% !TeX program = xelatex
% 在本目录编译。正文为保留原公式、图表和字形的矢量 PDF，不是可逐字编辑的 LaTeX。
\ifdefined\HandoutMode
  \documentclass[aspectratio=169,10pt,handout]{beamer}
\else
  \documentclass[aspectratio=169,10pt]{beamer}
\fi
\usepackage[UTF8,fontset=fandol]{ctex}
\IfFontExistsTF{Microsoft YaHei}{\setCJKsansfont{Microsoft YaHei}}{}
\usetheme{GaoZhongYuan}
\title{3.1.2 表示函数的方法}
\author{}
\date{}
\begin{document}
\title{3.1.2 表示函数的方法}\campusmaketitle

% 原 PPT 第 2 页；共 9 个显示状态。
\begin{frame}<1-9>[t]{复习导入}
\begin{campuscanvas}
\only<1|handout:0>{\node[anchor=north west,inner sep=0] at (25.600,116) {\includegraphics[page=2,width=153.6000mm]{source-content/content.pdf}};}
\only<2|handout:0>{\node[anchor=north west,inner sep=0] at (25.600,116) {\includegraphics[page=3,width=153.6000mm]{source-content/content.pdf}};}
\only<3|handout:0>{\node[anchor=north west,inner sep=0] at (25.600,116) {\includegraphics[page=4,width=153.6000mm]{source-content/content.pdf}};}
\only<4|handout:0>{\node[anchor=north west,inner sep=0] at (25.600,116) {\includegraphics[page=5,width=153.6000mm]{source-content/content.pdf}};}
\only<5|handout:0>{\node[anchor=north west,inner sep=0] at (25.600,116) {\includegraphics[page=6,width=153.6000mm]{source-content/content.pdf}};}
\only<6|handout:0>{\node[anchor=north west,inner sep=0] at (25.600,116) {\includegraphics[page=7,width=153.6000mm]{source-content/content.pdf}};}
\only<7|handout:0>{\node[anchor=north west,inner sep=0] at (25.600,116) {\includegraphics[page=8,width=153.6000mm]{source-content/content.pdf}};}
\only<8|handout:0>{\node[anchor=north west,inner sep=0] at (25.600,116) {\includegraphics[page=9,width=153.6000mm]{source-content/content.pdf}};}
\only<9|handout:1>{\node[anchor=north west,inner sep=0] at (25.600,116) {\includegraphics[page=10,width=153.6000mm]{source-content/content.pdf}};}
\end{campuscanvas}
\end{frame}

% 原 PPT 第 3 页；共 4 个显示状态。
\begin{frame}<1-4>[t]{复习导入}
\begin{campuscanvas}
\only<1|handout:0>{\node[anchor=north west,inner sep=0] at (25.600,116) {\includegraphics[page=11,width=153.6000mm]{source-content/content.pdf}};}
\only<2|handout:0>{\node[anchor=north west,inner sep=0] at (25.600,116) {\includegraphics[page=12,width=153.6000mm]{source-content/content.pdf}};}
\only<3|handout:0>{\node[anchor=north west,inner sep=0] at (25.600,116) {\includegraphics[page=13,width=153.6000mm]{source-content/content.pdf}};}
\only<4|handout:1>{\node[anchor=north west,inner sep=0] at (25.600,116) {\includegraphics[page=14,width=153.6000mm]{source-content/content.pdf}};}
\end{campuscanvas}
\end{frame}

% 原 PPT 第 4 页；共 5 个显示状态。
\begin{frame}<1-5>[t]{新课学习}
\begin{campuscanvas}
\only<1|handout:0>{\node[anchor=north west,inner sep=0] at (25.600,116) {\includegraphics[page=15,width=153.6000mm]{source-content/content.pdf}};}
\only<2|handout:0>{\node[anchor=north west,inner sep=0] at (25.600,116) {\includegraphics[page=16,width=153.6000mm]{source-content/content.pdf}};}
\only<3|handout:0>{\node[anchor=north west,inner sep=0] at (25.600,116) {\includegraphics[page=17,width=153.6000mm]{source-content/content.pdf}};}
\only<4|handout:0>{\node[anchor=north west,inner sep=0] at (25.600,116) {\includegraphics[page=18,width=153.6000mm]{source-content/content.pdf}};}
\only<5|handout:1>{\node[anchor=north west,inner sep=0] at (25.600,116) {\includegraphics[page=19,width=153.6000mm]{source-content/content.pdf}};}
\end{campuscanvas}
\end{frame}

% 原 PPT 第 5 页；共 10 个显示状态。
\begin{frame}<1-10>[t]{例题解析}
\begin{campuscanvas}
\only<1|handout:0>{\node[anchor=north west,inner sep=0] at (47.660,116) {\includegraphics[page=20,width=148.0851mm]{source-content/content.pdf}};}
\only<2|handout:0>{\node[anchor=north west,inner sep=0] at (47.660,116) {\includegraphics[page=21,width=148.0851mm]{source-content/content.pdf}};}
\only<3|handout:0>{\node[anchor=north west,inner sep=0] at (47.660,116) {\includegraphics[page=22,width=148.0851mm]{source-content/content.pdf}};}
\only<4|handout:0>{\node[anchor=north west,inner sep=0] at (47.660,116) {\includegraphics[page=23,width=148.0851mm]{source-content/content.pdf}};}
\only<5|handout:0>{\node[anchor=north west,inner sep=0] at (47.660,116) {\includegraphics[page=24,width=148.0851mm]{source-content/content.pdf}};}
\only<6|handout:0>{\node[anchor=north west,inner sep=0] at (47.660,116) {\includegraphics[page=25,width=148.0851mm]{source-content/content.pdf}};}
\only<7|handout:0>{\node[anchor=north west,inner sep=0] at (47.660,116) {\includegraphics[page=26,width=148.0851mm]{source-content/content.pdf}};}
\only<8|handout:0>{\node[anchor=north west,inner sep=0] at (47.660,116) {\includegraphics[page=27,width=148.0851mm]{source-content/content.pdf}};}
\only<9|handout:0>{\node[anchor=north west,inner sep=0] at (47.660,116) {\includegraphics[page=28,width=148.0851mm]{source-content/content.pdf}};}
\only<10|handout:1>{\node[anchor=north west,inner sep=0] at (47.660,116) {\includegraphics[page=29,width=148.0851mm]{source-content/content.pdf}};}
\end{campuscanvas}
\end{frame}

% 原 PPT 第 6 页；共 6 个显示状态。
\begin{frame}<1-6>[t]{例题解析}
\begin{campuscanvas}
\only<1|handout:0>{\node[anchor=north west,inner sep=0] at (25.600,116) {\includegraphics[page=30,width=153.6000mm]{source-content/content.pdf}};}
\only<2|handout:0>{\node[anchor=north west,inner sep=0] at (25.600,116) {\includegraphics[page=31,width=153.6000mm]{source-content/content.pdf}};}
\only<3|handout:0>{\node[anchor=north west,inner sep=0] at (25.600,116) {\includegraphics[page=32,width=153.6000mm]{source-content/content.pdf}};}
\only<4|handout:0>{\node[anchor=north west,inner sep=0] at (25.600,116) {\includegraphics[page=33,width=153.6000mm]{source-content/content.pdf}};}
\only<5|handout:0>{\node[anchor=north west,inner sep=0] at (25.600,116) {\includegraphics[page=34,width=153.6000mm]{source-content/content.pdf}};}
\only<6|handout:1>{\node[anchor=north west,inner sep=0] at (25.600,116) {\includegraphics[page=35,width=153.6000mm]{source-content/content.pdf}};}
\end{campuscanvas}
\end{frame}

% 原 PPT 第 7 页；共 4 个显示状态。
\begin{frame}<1-4>[t]{新课学习}
\begin{campuscanvas}
\only<1|handout:0>{\node[anchor=north west,inner sep=0] at (25.600,116) {\includegraphics[page=36,width=153.6000mm]{source-content/content.pdf}};}
\only<2|handout:0>{\node[anchor=north west,inner sep=0] at (25.600,116) {\includegraphics[page=37,width=153.6000mm]{source-content/content.pdf}};}
\only<3|handout:0>{\node[anchor=north west,inner sep=0] at (25.600,116) {\includegraphics[page=38,width=153.6000mm]{source-content/content.pdf}};}
\only<4|handout:1>{\node[anchor=north west,inner sep=0] at (25.600,116) {\includegraphics[page=39,width=153.6000mm]{source-content/content.pdf}};}
\end{campuscanvas}
\end{frame}

% 原 PPT 第 8 页；共 4 个显示状态。
\begin{frame}<1-4>[t]{新课学习}
\begin{campuscanvas}
\only<1|handout:0>{\node[anchor=north west,inner sep=0] at (25.600,116) {\includegraphics[page=40,width=153.6000mm]{source-content/content.pdf}};}
\only<2|handout:0>{\node[anchor=north west,inner sep=0] at (25.600,116) {\includegraphics[page=41,width=153.6000mm]{source-content/content.pdf}};}
\only<3|handout:0>{\node[anchor=north west,inner sep=0] at (25.600,116) {\includegraphics[page=42,width=153.6000mm]{source-content/content.pdf}};}
\only<4|handout:1>{\node[anchor=north west,inner sep=0] at (25.600,116) {\includegraphics[page=43,width=153.6000mm]{source-content/content.pdf}};}
\end{campuscanvas}
\end{frame}

% 原 PPT 第 9 页；共 7 个显示状态。
\begin{frame}<1-7>[t]{新课学习}
\begin{campuscanvas}
\only<1|handout:0>{\node[anchor=north west,inner sep=0] at (25.600,116) {\includegraphics[page=44,width=153.6000mm]{source-content/content.pdf}};}
\only<2|handout:0>{\node[anchor=north west,inner sep=0] at (25.600,116) {\includegraphics[page=45,width=153.6000mm]{source-content/content.pdf}};}
\only<3|handout:0>{\node[anchor=north west,inner sep=0] at (25.600,116) {\includegraphics[page=46,width=153.6000mm]{source-content/content.pdf}};}
\only<4|handout:0>{\node[anchor=north west,inner sep=0] at (25.600,116) {\includegraphics[page=47,width=153.6000mm]{source-content/content.pdf}};}
\only<5|handout:0>{\node[anchor=north west,inner sep=0] at (25.600,116) {\includegraphics[page=48,width=153.6000mm]{source-content/content.pdf}};}
\only<6|handout:0>{\node[anchor=north west,inner sep=0] at (25.600,116) {\includegraphics[page=49,width=153.6000mm]{source-content/content.pdf}};}
\only<7|handout:1>{\node[anchor=north west,inner sep=0] at (25.600,116) {\includegraphics[page=50,width=153.6000mm]{source-content/content.pdf}};}
\end{campuscanvas}
\end{frame}

% 原 PPT 第 10 页；共 2 个显示状态。
\begin{frame}<1-2>[t]{做一做}
\begin{campuscanvas}
\only<1|handout:0>{\node[anchor=north west,inner sep=0] at (25.600,116) {\includegraphics[page=51,width=153.6000mm]{source-content/content.pdf}};}
\only<2|handout:1>{\node[anchor=north west,inner sep=0] at (25.600,116) {\includegraphics[page=52,width=153.6000mm]{source-content/content.pdf}};}
\end{campuscanvas}
\end{frame}

% 原 PPT 第 11 页；共 2 个显示状态。
\begin{frame}<1-2>[t]{新课学习}
\begin{campuscanvas}
\only<1|handout:0>{\node[anchor=north west,inner sep=0] at (41.290,116) {\includegraphics[page=53,width=149.6774mm]{source-content/content.pdf}};}
\only<2|handout:1>{\node[anchor=north west,inner sep=0] at (41.290,116) {\includegraphics[page=54,width=149.6774mm]{source-content/content.pdf}};}
\end{campuscanvas}
\end{frame}

% 原 PPT 第 12 页；共 5 个显示状态。
\begin{frame}<1-5>[t]{新课学习}
\begin{campuscanvas}
\only<1|handout:0>{\node[anchor=north west,inner sep=0] at (25.600,116) {\includegraphics[page=55,width=153.6000mm]{source-content/content.pdf}};}
\only<2|handout:0>{\node[anchor=north west,inner sep=0] at (25.600,116) {\includegraphics[page=56,width=153.6000mm]{source-content/content.pdf}};}
\only<3|handout:0>{\node[anchor=north west,inner sep=0] at (25.600,116) {\includegraphics[page=57,width=153.6000mm]{source-content/content.pdf}};}
\only<4|handout:0>{\node[anchor=north west,inner sep=0] at (25.600,116) {\includegraphics[page=58,width=153.6000mm]{source-content/content.pdf}};}
\only<5|handout:1>{\node[anchor=north west,inner sep=0] at (25.600,116) {\includegraphics[page=59,width=153.6000mm]{source-content/content.pdf}};}
\end{campuscanvas}
\end{frame}

% 原 PPT 第 13 页；共 5 个显示状态。
\begin{frame}<1-5>[t]{例题解析}
\begin{campuscanvas}
\only<1|handout:0>{\node[anchor=north west,inner sep=0] at (28.132,116) {\includegraphics[page=60,width=152.9670mm]{source-content/content.pdf}};}
\only<2|handout:0>{\node[anchor=north west,inner sep=0] at (28.132,116) {\includegraphics[page=61,width=152.9670mm]{source-content/content.pdf}};}
\only<3|handout:0>{\node[anchor=north west,inner sep=0] at (28.132,116) {\includegraphics[page=62,width=152.9670mm]{source-content/content.pdf}};}
\only<4|handout:0>{\node[anchor=north west,inner sep=0] at (28.132,116) {\includegraphics[page=63,width=152.9670mm]{source-content/content.pdf}};}
\only<5|handout:1>{\node[anchor=north west,inner sep=0] at (28.132,116) {\includegraphics[page=64,width=152.9670mm]{source-content/content.pdf}};}
\end{campuscanvas}
\end{frame}

% 原 PPT 第 14 页；共 3 个显示状态。
\begin{frame}<1-3>[t]{做一做}
\begin{campuscanvas}
\only<1|handout:0>{\node[anchor=north west,inner sep=0] at (25.600,116) {\includegraphics[page=65,width=153.6000mm]{source-content/content.pdf}};}
\only<2|handout:0>{\node[anchor=north west,inner sep=0] at (25.600,116) {\includegraphics[page=66,width=153.6000mm]{source-content/content.pdf}};}
\only<3|handout:1>{\node[anchor=north west,inner sep=0] at (25.600,116) {\includegraphics[page=67,width=153.6000mm]{source-content/content.pdf}};}
\end{campuscanvas}
\end{frame}

% 原 PPT 第 15 页；共 4 个显示状态。
\begin{frame}<1-4>[t]{做一做}
\begin{campuscanvas}
\only<1|handout:0>{\node[anchor=north west,inner sep=0] at (25.600,116) {\includegraphics[page=68,width=153.6000mm]{source-content/content.pdf}};}
\only<2|handout:0>{\node[anchor=north west,inner sep=0] at (25.600,116) {\includegraphics[page=69,width=153.6000mm]{source-content/content.pdf}};}
\only<3|handout:0>{\node[anchor=north west,inner sep=0] at (25.600,116) {\includegraphics[page=70,width=153.6000mm]{source-content/content.pdf}};}
\only<4|handout:1>{\node[anchor=north west,inner sep=0] at (25.600,116) {\includegraphics[page=71,width=153.6000mm]{source-content/content.pdf}};}
\end{campuscanvas}
\end{frame}

% 原 PPT 第 16 页；共 1 个显示状态。
\begin{frame}<1-1>[t]{归纳小结}
\begin{campuscanvas}
\only<1|handout:1>{\node[anchor=north west,inner sep=0] at (25.600,116) {\includegraphics[page=72,width=153.6000mm]{source-content/content.pdf}};}
\end{campuscanvas}
\end{frame}

% 原 PPT 第 17 页；共 4 个显示状态。
\begin{frame}<1-4>[t]{归纳小结}
\begin{campuscanvas}
\only<1|handout:0>{\node[anchor=north west,inner sep=0] at (25.600,116) {\includegraphics[page=73,width=153.6000mm]{source-content/content.pdf}};}
\only<2|handout:0>{\node[anchor=north west,inner sep=0] at (25.600,116) {\includegraphics[page=74,width=153.6000mm]{source-content/content.pdf}};}
\only<3|handout:0>{\node[anchor=north west,inner sep=0] at (25.600,116) {\includegraphics[page=75,width=153.6000mm]{source-content/content.pdf}};}
\only<4|handout:1>{\node[anchor=north west,inner sep=0] at (25.600,116) {\includegraphics[page=76,width=153.6000mm]{source-content/content.pdf}};}
\end{campuscanvas}
\end{frame}

% 原 PPT 第 18 页；共 6 个显示状态。
\begin{frame}<1-6>[t]{课堂练习}
\begin{campuscanvas}
\only<1|handout:0>{\node[anchor=north west,inner sep=0] at (25.600,116) {\includegraphics[page=77,width=153.6000mm]{source-content/content.pdf}};}
\only<2|handout:0>{\node[anchor=north west,inner sep=0] at (25.600,116) {\includegraphics[page=78,width=153.6000mm]{source-content/content.pdf}};}
\only<3|handout:0>{\node[anchor=north west,inner sep=0] at (25.600,116) {\includegraphics[page=79,width=153.6000mm]{source-content/content.pdf}};}
\only<4|handout:0>{\node[anchor=north west,inner sep=0] at (25.600,116) {\includegraphics[page=80,width=153.6000mm]{source-content/content.pdf}};}
\only<5|handout:0>{\node[anchor=north west,inner sep=0] at (25.600,116) {\includegraphics[page=81,width=153.6000mm]{source-content/content.pdf}};}
\only<6|handout:1>{\node[anchor=north west,inner sep=0] at (25.600,116) {\includegraphics[page=82,width=153.6000mm]{source-content/content.pdf}};}
\end{campuscanvas}
\end{frame}

% 原 PPT 第 19 页；共 7 个显示状态。
\begin{frame}<1-7>[t]{课堂练习}
\begin{campuscanvas}
\only<1|handout:0>{\node[anchor=north west,inner sep=0] at (50.169,116) {\includegraphics[page=83,width=147.4576mm]{source-content/content.pdf}};}
\only<2|handout:0>{\node[anchor=north west,inner sep=0] at (50.169,116) {\includegraphics[page=84,width=147.4576mm]{source-content/content.pdf}};}
\only<3|handout:0>{\node[anchor=north west,inner sep=0] at (50.169,116) {\includegraphics[page=85,width=147.4576mm]{source-content/content.pdf}};}
\only<4|handout:0>{\node[anchor=north west,inner sep=0] at (50.169,116) {\includegraphics[page=86,width=147.4576mm]{source-content/content.pdf}};}
\only<5|handout:0>{\node[anchor=north west,inner sep=0] at (50.169,116) {\includegraphics[page=87,width=147.4576mm]{source-content/content.pdf}};}
\only<6|handout:0>{\node[anchor=north west,inner sep=0] at (50.169,116) {\includegraphics[page=88,width=147.4576mm]{source-content/content.pdf}};}
\only<7|handout:1>{\node[anchor=north west,inner sep=0] at (50.169,116) {\includegraphics[page=89,width=147.4576mm]{source-content/content.pdf}};}
\end{campuscanvas}
\end{frame}

% 原 PPT 第 20 页；共 8 个显示状态。
\begin{frame}<1-8>[t]{课堂练习}
\begin{campuscanvas}
\only<1|handout:0>{\node[anchor=north west,inner sep=0] at (25.600,116) {\includegraphics[page=90,width=153.6000mm]{source-content/content.pdf}};}
\only<2|handout:0>{\node[anchor=north west,inner sep=0] at (25.600,116) {\includegraphics[page=91,width=153.6000mm]{source-content/content.pdf}};}
\only<3|handout:0>{\node[anchor=north west,inner sep=0] at (25.600,116) {\includegraphics[page=92,width=153.6000mm]{source-content/content.pdf}};}
\only<4|handout:0>{\node[anchor=north west,inner sep=0] at (25.600,116) {\includegraphics[page=93,width=153.6000mm]{source-content/content.pdf}};}
\only<5|handout:0>{\node[anchor=north west,inner sep=0] at (25.600,116) {\includegraphics[page=94,width=153.6000mm]{source-content/content.pdf}};}
\only<6|handout:0>{\node[anchor=north west,inner sep=0] at (25.600,116) {\includegraphics[page=95,width=153.6000mm]{source-content/content.pdf}};}
\only<7|handout:0>{\node[anchor=north west,inner sep=0] at (25.600,116) {\includegraphics[page=96,width=153.6000mm]{source-content/content.pdf}};}
\only<8|handout:1>{\node[anchor=north west,inner sep=0] at (25.600,116) {\includegraphics[page=97,width=153.6000mm]{source-content/content.pdf}};}
\end{campuscanvas}
\end{frame}

% 原 PPT 第 21 页；共 1 个显示状态。
\begin{frame}<1-1>[t]{方法提炼}
\begin{campuscanvas}
\only<1|handout:1>{\node[anchor=north west,inner sep=0] at (25.600,116) {\includegraphics[page=98,width=153.6000mm]{source-content/content.pdf}};}
\end{campuscanvas}
\end{frame}

% 原 PPT 第 22 页；共 3 个显示状态。
\begin{frame}<1-3>[t]{课堂总结}
\begin{campuscanvas}
\only<1|handout:0>{\node[anchor=north west,inner sep=0] at (25.600,116) {\includegraphics[page=99,width=153.6000mm]{source-content/content.pdf}};}
\only<2|handout:0>{\node[anchor=north west,inner sep=0] at (25.600,116) {\includegraphics[page=100,width=153.6000mm]{source-content/content.pdf}};}
\only<3|handout:1>{\node[anchor=north west,inner sep=0] at (25.600,116) {\includegraphics[page=101,width=153.6000mm]{source-content/content.pdf}};}
\end{campuscanvas}
\end{frame}

\title{专题：函数解析式的求法}\campusmaketitle

% 原 PPT 第 24 页；共 8 个显示状态。
\begin{frame}<1-8>[t]{\resizebox{94mm}{!}{题型一：求函数的解析式——代入法}}
\begin{campuscanvas}
\only<1|handout:0>{\node[anchor=north west,inner sep=0] at (25.600,116) {\includegraphics[page=103,width=153.6000mm]{source-content/content.pdf}};}
\only<2|handout:0>{\node[anchor=north west,inner sep=0] at (25.600,116) {\includegraphics[page=104,width=153.6000mm]{source-content/content.pdf}};}
\only<3|handout:0>{\node[anchor=north west,inner sep=0] at (25.600,116) {\includegraphics[page=105,width=153.6000mm]{source-content/content.pdf}};}
\only<4|handout:0>{\node[anchor=north west,inner sep=0] at (25.600,116) {\includegraphics[page=106,width=153.6000mm]{source-content/content.pdf}};}
\only<5|handout:0>{\node[anchor=north west,inner sep=0] at (25.600,116) {\includegraphics[page=107,width=153.6000mm]{source-content/content.pdf}};}
\only<6|handout:0>{\node[anchor=north west,inner sep=0] at (25.600,116) {\includegraphics[page=108,width=153.6000mm]{source-content/content.pdf}};}
\only<7|handout:0>{\node[anchor=north west,inner sep=0] at (25.600,116) {\includegraphics[page=109,width=153.6000mm]{source-content/content.pdf}};}
\only<8|handout:1>{\node[anchor=north west,inner sep=0] at (25.600,116) {\includegraphics[page=110,width=153.6000mm]{source-content/content.pdf}};}
\end{campuscanvas}
\end{frame}

% 原 PPT 第 25 页；共 5 个显示状态。
\begin{frame}<1-5>[t]{\resizebox{94mm}{!}{题型一：求函数的解析式——代入法}}
\begin{campuscanvas}
\only<1|handout:0>{\node[anchor=north west,inner sep=0] at (25.600,116) {\includegraphics[page=111,width=153.6000mm]{source-content/content.pdf}};}
\only<2|handout:0>{\node[anchor=north west,inner sep=0] at (25.600,116) {\includegraphics[page=112,width=153.6000mm]{source-content/content.pdf}};}
\only<3|handout:0>{\node[anchor=north west,inner sep=0] at (25.600,116) {\includegraphics[page=113,width=153.6000mm]{source-content/content.pdf}};}
\only<4|handout:0>{\node[anchor=north west,inner sep=0] at (25.600,116) {\includegraphics[page=114,width=153.6000mm]{source-content/content.pdf}};}
\only<5|handout:1>{\node[anchor=north west,inner sep=0] at (25.600,116) {\includegraphics[page=115,width=153.6000mm]{source-content/content.pdf}};}
\end{campuscanvas}
\end{frame}

% 原 PPT 第 26 页；共 8 个显示状态。
\begin{frame}<1-8>[t]{\resizebox{94mm}{!}{题型二：求函数的解析式——换元法}}
\begin{campuscanvas}
\only<1|handout:0>{\node[anchor=north west,inner sep=0] at (25.600,116) {\includegraphics[page=116,width=153.6000mm]{source-content/content.pdf}};}
\only<2|handout:0>{\node[anchor=north west,inner sep=0] at (25.600,116) {\includegraphics[page=117,width=153.6000mm]{source-content/content.pdf}};}
\only<3|handout:0>{\node[anchor=north west,inner sep=0] at (25.600,116) {\includegraphics[page=118,width=153.6000mm]{source-content/content.pdf}};}
\only<4|handout:0>{\node[anchor=north west,inner sep=0] at (25.600,116) {\includegraphics[page=119,width=153.6000mm]{source-content/content.pdf}};}
\only<5|handout:0>{\node[anchor=north west,inner sep=0] at (25.600,116) {\includegraphics[page=120,width=153.6000mm]{source-content/content.pdf}};}
\only<6|handout:0>{\node[anchor=north west,inner sep=0] at (25.600,116) {\includegraphics[page=121,width=153.6000mm]{source-content/content.pdf}};}
\only<7|handout:0>{\node[anchor=north west,inner sep=0] at (25.600,116) {\includegraphics[page=122,width=153.6000mm]{source-content/content.pdf}};}
\only<8|handout:1>{\node[anchor=north west,inner sep=0] at (25.600,116) {\includegraphics[page=123,width=153.6000mm]{source-content/content.pdf}};}
\end{campuscanvas}
\end{frame}

% 原 PPT 第 27 页；共 3 个显示状态。
\begin{frame}<1-3>[t]{\resizebox{94mm}{!}{题型二：求函数的解析式——换元法}}
\begin{campuscanvas}
\only<1|handout:0>{\node[anchor=north west,inner sep=0] at (25.600,116) {\includegraphics[page=124,width=153.6000mm]{source-content/content.pdf}};}
\only<2|handout:0>{\node[anchor=north west,inner sep=0] at (25.600,116) {\includegraphics[page=125,width=153.6000mm]{source-content/content.pdf}};}
\only<3|handout:1>{\node[anchor=north west,inner sep=0] at (25.600,116) {\includegraphics[page=126,width=153.6000mm]{source-content/content.pdf}};}
\end{campuscanvas}
\end{frame}

% 原 PPT 第 28 页；共 7 个显示状态。
\begin{frame}<1-7>[t]{\resizebox{94mm}{!}{题型三：求函数的解析式——配凑法}}
\begin{campuscanvas}
\only<1|handout:0>{\node[anchor=north west,inner sep=0] at (25.600,116) {\includegraphics[page=127,width=153.6000mm]{source-content/content.pdf}};}
\only<2|handout:0>{\node[anchor=north west,inner sep=0] at (25.600,116) {\includegraphics[page=128,width=153.6000mm]{source-content/content.pdf}};}
\only<3|handout:0>{\node[anchor=north west,inner sep=0] at (25.600,116) {\includegraphics[page=129,width=153.6000mm]{source-content/content.pdf}};}
\only<4|handout:0>{\node[anchor=north west,inner sep=0] at (25.600,116) {\includegraphics[page=130,width=153.6000mm]{source-content/content.pdf}};}
\only<5|handout:0>{\node[anchor=north west,inner sep=0] at (25.600,116) {\includegraphics[page=131,width=153.6000mm]{source-content/content.pdf}};}
\only<6|handout:0>{\node[anchor=north west,inner sep=0] at (25.600,116) {\includegraphics[page=132,width=153.6000mm]{source-content/content.pdf}};}
\only<7|handout:1>{\node[anchor=north west,inner sep=0] at (25.600,116) {\includegraphics[page=133,width=153.6000mm]{source-content/content.pdf}};}
\end{campuscanvas}
\end{frame}

% 原 PPT 第 29 页；共 4 个显示状态。
\begin{frame}<1-4>[t]{\resizebox{94mm}{!}{题型三：求函数的解析式——配凑法}}
\begin{campuscanvas}
\only<1|handout:0>{\node[anchor=north west,inner sep=0] at (25.600,116) {\includegraphics[page=134,width=153.6000mm]{source-content/content.pdf}};}
\only<2|handout:0>{\node[anchor=north west,inner sep=0] at (25.600,116) {\includegraphics[page=135,width=153.6000mm]{source-content/content.pdf}};}
\only<3|handout:0>{\node[anchor=north west,inner sep=0] at (25.600,116) {\includegraphics[page=136,width=153.6000mm]{source-content/content.pdf}};}
\only<4|handout:1>{\node[anchor=north west,inner sep=0] at (25.600,116) {\includegraphics[page=137,width=153.6000mm]{source-content/content.pdf}};}
\end{campuscanvas}
\end{frame}

% 原 PPT 第 30 页；共 3 个显示状态。
\begin{frame}<1-3>[t]{\resizebox{94mm}{!}{题型四：求函数的解析式——待定系数法}}
\begin{campuscanvas}
\only<1|handout:0>{\node[anchor=north west,inner sep=0] at (25.600,116) {\includegraphics[page=138,width=153.6000mm]{source-content/content.pdf}};}
\only<2|handout:0>{\node[anchor=north west,inner sep=0] at (25.600,116) {\includegraphics[page=139,width=153.6000mm]{source-content/content.pdf}};}
\only<3|handout:1>{\node[anchor=north west,inner sep=0] at (25.600,116) {\includegraphics[page=140,width=153.6000mm]{source-content/content.pdf}};}
\end{campuscanvas}
\end{frame}

% 原 PPT 第 31 页；共 4 个显示状态。
\begin{frame}<1-4>[t]{\resizebox{94mm}{!}{题型四：求函数的解析式——待定系数法}}
\begin{campuscanvas}
\only<1|handout:0>{\node[anchor=north west,inner sep=0] at (25.600,116) {\includegraphics[page=141,width=153.6000mm]{source-content/content.pdf}};}
\only<2|handout:0>{\node[anchor=north west,inner sep=0] at (25.600,116) {\includegraphics[page=142,width=153.6000mm]{source-content/content.pdf}};}
\only<3|handout:0>{\node[anchor=north west,inner sep=0] at (25.600,116) {\includegraphics[page=143,width=153.6000mm]{source-content/content.pdf}};}
\only<4|handout:1>{\node[anchor=north west,inner sep=0] at (25.600,116) {\includegraphics[page=144,width=153.6000mm]{source-content/content.pdf}};}
\end{campuscanvas}
\end{frame}

% 原 PPT 第 32 页；共 4 个显示状态。
\begin{frame}<1-4>[t]{\resizebox{94mm}{!}{题型四：求函数的解析式——待定系数法}}
\begin{campuscanvas}
\only<1|handout:0>{\node[anchor=north west,inner sep=0] at (25.600,116) {\includegraphics[page=145,width=153.6000mm]{source-content/content.pdf}};}
\only<2|handout:0>{\node[anchor=north west,inner sep=0] at (25.600,116) {\includegraphics[page=146,width=153.6000mm]{source-content/content.pdf}};}
\only<3|handout:0>{\node[anchor=north west,inner sep=0] at (25.600,116) {\includegraphics[page=147,width=153.6000mm]{source-content/content.pdf}};}
\only<4|handout:1>{\node[anchor=north west,inner sep=0] at (25.600,116) {\includegraphics[page=148,width=153.6000mm]{source-content/content.pdf}};}
\end{campuscanvas}
\end{frame}

% 原 PPT 第 33 页；共 3 个显示状态。
\begin{frame}<1-3>[t]{\resizebox{94mm}{!}{题型五：求函数的解析式——方程组法}}
\begin{campuscanvas}
\only<1|handout:0>{\node[anchor=north west,inner sep=0] at (25.600,116) {\includegraphics[page=149,width=153.6000mm]{source-content/content.pdf}};}
\only<2|handout:0>{\node[anchor=north west,inner sep=0] at (25.600,116) {\includegraphics[page=150,width=153.6000mm]{source-content/content.pdf}};}
\only<3|handout:1>{\node[anchor=north west,inner sep=0] at (25.600,116) {\includegraphics[page=151,width=153.6000mm]{source-content/content.pdf}};}
\end{campuscanvas}
\end{frame}

% 原 PPT 第 34 页；共 2 个显示状态。
\begin{frame}<1-2>[t]{\resizebox{94mm}{!}{题型五：求函数的解析式——方程组法}}
\begin{campuscanvas}
\only<1|handout:0>{\node[anchor=north west,inner sep=0] at (25.600,116) {\includegraphics[page=152,width=153.6000mm]{source-content/content.pdf}};}
\only<2|handout:1>{\node[anchor=north west,inner sep=0] at (25.600,116) {\includegraphics[page=153,width=153.6000mm]{source-content/content.pdf}};}
\end{campuscanvas}
\end{frame}

% 原 PPT 第 35 页；共 2 个显示状态。
\begin{frame}<1-2>[t]{\resizebox{94mm}{!}{题型五：求函数的解析式——方程组法}}
\begin{campuscanvas}
\only<1|handout:0>{\node[anchor=north west,inner sep=0] at (25.600,116) {\includegraphics[page=154,width=153.6000mm]{source-content/content.pdf}};}
\only<2|handout:1>{\node[anchor=north west,inner sep=0] at (25.600,116) {\includegraphics[page=155,width=153.6000mm]{source-content/content.pdf}};}
\end{campuscanvas}
\end{frame}

% 原 PPT 第 36 页；共 1 个显示状态。
\begin{frame}<1-1>[t]{课堂小结}
\begin{campuscanvas}
\only<1|handout:1>{\node[anchor=north west,inner sep=0] at (25.600,116) {\includegraphics[page=156,width=153.6000mm]{source-content/content.pdf}};}
\end{campuscanvas}
\end{frame}

% 原 PPT 第 37 页；共 1 个显示状态。
\begin{frame}<1-1>[t]{课堂小结}
\begin{campuscanvas}
\only<1|handout:1>{\node[anchor=north west,inner sep=0] at (25.600,116) {\includegraphics[page=157,width=153.6000mm]{source-content/content.pdf}};}
\end{campuscanvas}
\end{frame}

% 原 PPT 第 38 页；共 1 个显示状态。
\begin{frame}<1-1>[t]{课堂小结}
\begin{campuscanvas}
\only<1|handout:1>{\node[anchor=north west,inner sep=0] at (25.600,116) {\includegraphics[page=158,width=153.6000mm]{source-content/content.pdf}};}
\end{campuscanvas}
\end{frame}

\end{document}
